\documentclass[../../../main]{subfiles}

\ifSubfilesClassLoaded{%
    \zexternaldocument*{../1.Category-Theory/exported/aux/main}%
    \zexternaldocument*{../1.Additive-and-abelian-categories/exported/aux/main}%
    \zexternaldocument*{../2.Complexes-and-cohomology/exported/aux/main}%
    \zexternaldocument*{../3.Cones-and-homotopy/exported/aux/main}%
    \zexternaldocument*{../4.Projectives-and-injectives/exported/aux/main}%
}{}

\begin{document}
\section{Delta functors}

\begin{definition} \label{def:delta-functor}
  Let \(\category{A}, \category{B}\) be two \hyperref[def:abelian-category]{abelian categories}.
  A \textbf{cohomological \(\delta\)-functor} is a collection of \hyperref[def:additive-functor]{additive functors} \(T^{n}: \category{A} \to \category{B}\) for \(n \in \mathbb{N}\) together with, for every \hyperref[def:short-exact-sequence]{short exact sequence}
  \begin{center}
    \begin{tikzcd}
      0 \arrow[r] & A \arrow[r] & B \arrow[r] & C \arrow[r] & 0
    \end{tikzcd}
  \end{center}
  a morphism \(\delta^{n}: T^{n}(C) \to T^{n+1}(A)\). We take the convention that \(\forall n < 0 \quad T^{n} = 0\).
  Furthermore
  \begin{enumerate}
    \item For each \hyperref[def:short-exact-sequence]{short exact sequence} as above there exists a \hyperref[def:exact-complex-sequence]{long exact sequence}
          \begin{center}
            \begin{tikzcd}
              \cdots \arrow[r] & T^{n-1}(B) \arrow[r] \arrow[d, phantom, ""{coordinate, name=X}] & T^{n-1}(C) \arrow[dll, rounded corners, to path={
        -- ([xshift=2ex]\tikztostart.east)
        |- (X) [near end]\tikztonodes
        -| ([xshift=-2ex]\tikztotarget.west)
        -- (\tikztotarget)}, "\delta^{n-1}" description] \\
      T^{n}(A) \arrow[r] & T^{n}(B) \arrow[r] \arrow[d, phantom, ""{coordinate, name=Y}] & T^{n}(C) \arrow[dll, rounded corners, to path={
        -- ([xshift=2ex]\tikztostart.east)
        |- (Y) [near end]\tikztonodes
        -| ([xshift=-2ex]\tikztotarget.west)
        -- (\tikztotarget)}, "\delta^{n}" description] \\
      T^{n+1}(A) \arrow[r] & T^{n+1}(B) \arrow[r] & \cdots
            \end{tikzcd}
          \end{center}
          In particular \(T^{0}\) is \hyperref[def:exact-complex-functor]{left exact}.
    \item For each morphism of \hyperref[def:short-exact-sequence]{short exact sequences}
          \begin{equation*}
            \left( 0 \to A' \to B' \to C' \to 0 \right) \mapsto \left( 0 \to A \to B \to C \to 0 \right)
          \end{equation*}
          the \(\delta^{n}\) give commutative diagrams
          \begin{center}
            \begin{tikzcd}
              T^{n}(C') \arrow[r, "\delta^{n}"] \arrow[d] & T^{n+1}(A') \arrow[d] \\
              T^{n}(C) \arrow[r, "\delta^{n}"] & T^{n+1}(A)
            \end{tikzcd}
          \end{center}
  \end{enumerate}
\end{definition}

\begin{example}
  \hyperref[def:cohomology]{Cohomology} is a \hyperref[def:delta-functor]{cohomological \(\delta\)-functor} \(\cohomology{\bullet}\) from \(\Complex[\geq 0]{\category{A}}\) to \(\category{A}\).
\end{example}

\begin{definition} \label{def:morphism-delta-functor}
  A morphism \(S \to T\) of \hyperref[def:delta-functor]{cohomological \(\delta\)-functors} is a collection of \hyperref[def:natural-transformation]{natural transformations} \(S^{n} \to T^{n}\) that commute with \(\delta\).
\end{definition}

\begin{definition} \label{def:universal-delta-functor}
  A \hyperref[def:delta-functor]{cohomological \(\delta\)-functor} is \textbf{universal} if given another \hyperref[def:delta-functor]{cohomological \(\delta\)-functor} \(S\) and \(f^{0}: T^{0} \to S^{0}\) there exists a unique \hyperref[def:morphism-delta-functor]{morphism of \(\delta\)-functors} \(T \to S\) extending \(f^{0}\).
\end{definition}


\end{document}
